\section{多元函数的极限与连续 II}

\subsection{Euclid 空间 continue}

\begin{definition}[基本列]
    设 $\{\vect{x}_n\}$ 是 $\mathbb{R}^{n}$ 中的点列，若：
    $$
    \forall\,\varepsilon > 0: \exists\,N \in \mathbb{N}^{+}: \forall\,i, j > N: \norm{\vect{x}_i - \vect{x}_j} < \varepsilon
    $$
    那么称 $\{\vect{x}_n\}$ 是一个基本列。
\end{definition}

\begin{theorem}
    基本列收敛。

    \begin{proof}
        对每一个分量应用基本列定理即可。
    \end{proof}
\end{theorem}

\subsubsection{点集的相关概念}

接下来讨论 $\mathbb{R}^{n}$ 中的特殊点集的定义和其具有的性质。

\begin{definition}[开集和闭集]
    设 $E \in \mathbb{R}^{n}$，若
    $$
    \forall\,\vect{x} \in E: \exists\,\varepsilon > 0: B_{\varepsilon}(\vect{x}) \subset E
    $$
    那么称 $E$ 为\textbf{开集}。

    若一个集合的补集是开集，则称该集合为\textbf{闭集}。
\end{definition}

\begin{remark}
    空集 $\varnothing$ 和全集 $\mathbb{R}^n$ 既是开集，也是闭集。
\end{remark}

\begin{property}
    \begin{itemize}
        \item 任意多个开集的并集是开集，有限多个开集的交集是开集（不能无限多个，反例：区间套）；
        \item 任意多个闭集的交集是闭集，有限多个闭集的并集是闭集（不能无限多个，反例：区间套）；
    \end{itemize}
\end{property}

\begin{definition}[内点、外点、边界点]
    设 $E \subset \mathbb{R}^{n}$，那么：
    \begin{itemize}
        \item 若点 $\vect{x} \in E$ 满足：
        $$
        \exists\,\varepsilon > 0: B_{\varepsilon}(\vect{x}) \subset E
        $$
        则称 $\vect{x}$ 是 $E$ 的一个\textbf{内点}。所有 $E$ 的内点组成的集合称为 $E$ 的\textbf{内部}，记作 $E^o$。
        
        \item 若点 $\vect{x}$ 是 $E^c$ 的内点，那么称 $\vect{x}$ 是 $E$ 的一个\textbf{外点}。所有 $E$ 的外点组成的集合称为 $E$ 的\textbf{外部}。
        
        \item 若点 $\vect{x}$ 满足：
        $$
        \forall\,\varepsilon > 0: \exists\,\vect{u}, \vect{v} \in B_{\varepsilon}(\vect{x}): \vect{u} \in E^o \land \vect{v} \in (E^c)^o
        $$
        那么称 $\vect{x}$ 是 $E$ 的一个\textbf{边界点}。所有 $E$ 的边界点组成的集合称为 $E$ 的\textbf{边界}，记作 $\partial E$。
    \end{itemize}
\end{definition}

\begin{definition}[聚点]
    设 $E \subseteq \mathbb{R}^{n}$，那么若 $\vect{a} \in \mathbb{R}^{n}$ 满足：
    $$
    \forall\, \varepsilon > 0: \ab(B_{\varepsilon}(\vect{a}) \setminus \{\vect{a}\}) \cap E \neq \varnothing
    $$
    那么称 $\vect{a}$ 是 $E$ 的一个\textbf{聚点}。
\end{definition}

\begin{remark}
    集合 $E$ 的聚点可以属于 $E$ 也可以不属于 $E$。若 $E$ 中某点不是 $E$ 的聚点，其又称为\textbf{孤立点}。
\end{remark}

\begin{definition}[导集]
    点集 $E$ 的所有聚点组成的集合称为 $E$ 的\textbf{导集}，记作 $E'$。记 $\overline{E} = E \cup E'$ 为 $E$ 的\textbf{闭包}。
\end{definition}

导集可以直观理解为，将 $E$ 的边界加上，将孤立点去除得到的点集。闭包可以理解为，将 $E$ 的边界全部加入 $E$ 形成一个闭集。

\begin{theorem}
    $E$ 是闭集等价于 $E$ 中的任意收敛点列极限均在 $E$ 中。
\end{theorem}

\begin{theorem}
    $E$ 的导集和闭包都是闭集。
\end{theorem}

下面讨论点集的连通性问题。首先定义连续映射的概念：

\begin{definition}
    给定映射
    $$
    \varphi = (\varphi_1(t), \varphi_2(t), \dots, \varphi_n(t)): [a, b] \to \mathbb{R}^{n}
    $$
    若所有分量函数 $\varphi_i(t)$ 均连续，则称 $\varphi$ 为\textbf{连续映射}，它的像称为一条\textbf{连续曲线}。
\end{definition}

然后我们就可定义道路连通和区域：

\begin{definition}[道路连通]
    设 $E \subset \mathbb{R}^{n}$，若对于任意 $\vect{x}, \vect{y} \in E$，存在连续映射
    $$
    \varphi: [0, 1] \to E,\ \varphi(0) = \vect{x} \land \varphi(1) = \vect{y}
    $$
    那么称 $E$ 是\textbf{道路连通}的。
\end{definition}

\begin{definition}[区域]
    对于一个 $\mathbb{R}^{n}$ 中的道路连通的开集，称为\textbf{（开）区域}。区域的闭包称为\textbf{闭区域}。
\end{definition}

\subsection{Euclid 空间的基本定理}

在实数集中的很多理论可以扩展导 Euclid 空间。

\begin{theorem}[闭集套定理]
    设 $\{E_k\}$ 是 $\mathbb{R}^{n}$ 上的非空闭集序列，满足 $E_1 \supset E_2 \supset \dots \supset E_n \supset \dots$，且 $\lim\limits_{k \to \infty} \operatorname{diam} E_k = 0$，则：
    $$ 
    \bigcap_{k=1}^{\infty} E_k = \{\vect{x}_0\}
    $$
    其中 $\operatorname{diam} E = \sup\limits_{\vect{x}, \vect{y} \in E} \{\norm{\vect{x} - \vect{y}}\}$ 称为 $E$ 的直径。 

    \begin{proof}
        可以直接对分量应用闭区间套定理。
    \end{proof}
\end{theorem}

\begin{theorem}[Bolzano-Weierstrass 定理]
    $\mathbb{R}^{n}$ 上的有界子列必有收敛子列。

    \begin{proof}
        依次对于每个分量，应用列紧性定理选择出一个子列即可。
    \end{proof}
\end{theorem}

\begin{definition}
    设 $S \subseteq \mathbb{R}^{n}$，若 $\mathbb{R}^{n}$ 的一组开集 $\{U_\alpha\}$ 满足 $S \subset \bigcup_{\alpha} U_\alpha$，则称 $\{U_\alpha\}$ 是 $S$ 的一个\textbf{开覆盖}。
    
    若从任意开覆盖中总能选出一个有限子覆盖，使得它们的并集仍包含 $S$，则称 $S$ 是\textbf{紧集}。
\end{definition}

对于禁集和有限闭集之间的关系，在 Euclid 空间中有定理：

\begin{theorem}
    设 $E \subset \mathbb{R}^n$，则下面几条等价：
    \begin{itemize}
        \item $E$ 是紧集；
        \item $E$ 是有界闭集；
        \item $E$ 中的任意无穷点列均有收敛子列，且极限均在 $E$ 中。
    \end{itemize}
\end{theorem}

\subsection{多元函数的极限与连续}

\begin{definition}[多元函数]
    设 $D \subset \mathbb{R}^n$，映射
    $$
    f: D \to \mathbb{R}
    $$
    称为 $n$ 元函数。
\end{definition}

\begin{definition}[多元函数极限]
    设 $D \subset \mathbb{R}^{n}, f: D \to \mathbb{R}$。$\vect{a} \in \mathbb{R}^{n}$ 是 $D$ 的一个聚点，$A \in \mathbb{R}$，那么若：
    $$
    \forall\,\varepsilon > 0: \exists\,\delta > 0: \forall\,\vect{x} \in D,\ \norm{\vect{x} - \vect{a}} < \delta: \abs{f(\vect{x}) - A} < \varepsilon
    $$
    就称函数 $f$ 在 $\vect{a}$ 上有极限 $A$，记做
    $$
    \lim_{\vect{x} \to \vect{a}} f(\vect{x}) = A
    $$
\end{definition}

一元函数中有连续的概念，在多元函数同样需要定义连续：

\begin{definition}[多元函数连续]
    设 $D \subset \mathbb{R}^{n}, f: D \to \mathbb{R}$。若 $\vect{a} \in D$ 是 $D$ 的一个内点，且：
    $$
    \forall\,\varepsilon > 0: \exists\,\delta > 0: \forall\,\vect{x} \in D,\ \norm{\vect{x} - \vect{a}} < \delta: \abs{f(\vect{x}) - f(\vect{a})} < \varepsilon
    $$
    即
    $$
    \lim_{\vect{x} \to \vect{a}} f(\vect{x}) = f(\vect{a})
    $$
    则称函数 $f$ 在 $\vect{a}$ 处连续。
\end{definition}